\section{Recoil molecular dynamics}\label{app:E.11}

A silicon atom struck with \SI{100}{eV} leaves its site at \SI{26}{km/s}, comes within about \SI{1}{\angstrom} of a neighbor, and within a few picoseconds has either fallen back into the lattice or left behind a vacancy and an interstitial that store several eV between them. Nothing in this sequence is a small perturbation of the harmonic crystal. Bonds break and re-form, and the outcome depends on the shape of the Born--Oppenheimer energy surface far from equilibrium. \Cref{sec:19.1} hands the recoil over to molecular dynamics on this energy surface, and \cref{sec:19.4,sec:19.6} need four products of the simulation: the displacement threshold $E_{\mathrm{d}}$ as a function of recoil direction (\cref{sec:7.4}), the distribution of the stored energy $E_{\mathrm{st}}$ at a given recoil energy $E_R$, the configurations of the defects that survive, and the barriers along the paths by which they recombine or move. Each comes from many trajectories on an energy surface fitted to density-functional theory, with a screened Coulomb repulsion for the close collisions, and from the statistics of those trajectories over recoil directions.

\subsection{Computed quantities}

Molecular dynamics integrates Newton's equations for the nuclei on a single energy surface,
\begin{equation}
  M_\kappa\,\ddot{\vect{R}}_I = -\frac{\partial E^{\mathrm{ad}}(\{\vect{R}\})}{\partial\vect{R}_I},
  \label{eq:E.11-newton}
\end{equation}
where $\vect{R}_I$ is the position of nucleus $I$ of species $\kappa$ and $E^{\mathrm{ad}}(\{\vect{R}\})$ is the electronic ground-state energy at fixed nuclear positions, the Born--Oppenheimer surface of \cref{sec:5.1}. A recoil simulation starts from a periodic crystal cell equilibrated at the base temperature and gives one atom (the primary knock-on atom) the kinetic energy $E_R$ in a chosen direction $\hat{\vect{r}}_R$. The simulation then integrates \cref{eq:E.11-newton} until the part of $E_R$ that is not stored has spread through the cell, and records what changed. The quantities read off a single trajectory are whether a defect survived, which defect, and the potential energy that the defect stores. Repeated over directions and energies, the trajectories give the probability $p_{\mathrm{FP}}(E_R, \hat{\vect{r}}_R)$ that a Frenkel pair survives, the threshold surface $E_{\mathrm{d}}(\hat{\vect{r}}_R)$, and the distribution of $E_{\mathrm{st}}$ that \cref{sec:19.4} bounds with calibration data.

Density-functional theory (\cref{app:E.1}) gives $E^{\mathrm{ad}}$ and its gradients at a cost that grows as the cube of the number of atoms. A plane-wave calculation with $N_{\mathrm{at}}$ atoms holds a number of occupied bands and a number of plane waves that are both proportional to $N_{\mathrm{at}}$, and orthogonalizing the bands costs their number squared times the basis size. A recoil cascade needs thousands of atoms, to keep the recoil away from its periodic images, and thousands of trajectories, to sample directions. Direct density-functional molecular dynamics can afford a few directions in a cell of a few hundred atoms, and serves as the reference; the production runs use an interatomic potential that reproduces $E^{\mathrm{ad}}$ at a cost linear in $N_{\mathrm{at}}$.

\subsection{Place on the hierarchy}

In the language of the truncation label $\mathfrak{T}$ of \cref{sec:8.5}, \cref{eq:E.11-newton} is the adiabatic limit, $n_B = 0$: the electrons stay in their ground state for every nuclear configuration and never take energy from the nuclei. It keeps the anharmonicity of $E^{\mathrm{ad}}$ to all orders, which the expansion of \cref{sec:5.7} cannot, and it lies beyond the impulse regime of the lattice label $\mathcal{K}$, since the struck atom collides with its neighbors as it leaves its site. Two approximations come with the adiabatic limit. The first is that the nuclei are classical, so zero-point motion is absent and the lattice starts at rest apart from thermal motion. For recoils above \SI{10}{eV} the missing zero-point motion does not matter for the mean threshold, since the kinetic energies exceed the vibrational energies of a few tens of meV by orders of magnitude~\cite[Sect.~II\,A]{sassi2022energy}. Zero-point motion does set the width of the threshold step, \SI{0.7}{eV} at \SI{12.5}{eV} (\cref{sec:19.1-handover}). The classical treatment of the nuclei also matters for the slow processes at millikelvin, where barriers are crossed by tunneling rather than by thermal activation (the defect paths below). The second approximation is that energy given to the electrons is lost from the description. At \SI{1}{keV} a silicon recoil already transfers part of its energy to electronic excitations, a channel that needs the nonadiabatic treatment of \cref{sec:19.2}; molecular dynamics then supplies the nuclear trajectory on which the electronic losses are evaluated.

\subsection{Interatomic potentials}

An interatomic potential replaces $E^{\mathrm{ad}}$ by a sum of local terms,
\begin{equation}
  E^{\mathrm{ad}}(\{\vect{R}\}) \approx \sum_I E^{\mathrm{at}}_\kappa(\mathcal{N}_I),
  \label{eq:E.11-local}
\end{equation}
where $\mathcal{N}_I$ is the set of positions of the neighbors of atom $I$ within a cutoff radius $r_{\mathrm{c}}$, measured from $\vect{R}_I$. Locality makes the cost linear in $N_{\mathrm{at}}$. It holds for covalent and metallic solids because the density matrix decays with distance, and it fails for long-range electrostatics in ionic materials, where the electrostatic energy is then added as an explicit Coulomb term.

The classical forms write $E^{\mathrm{at}}_\kappa$ as a handful of terms with physically motivated shapes. The Stillinger--Weber potential for silicon is a pair term plus a three-body term that penalizes deviations from the tetrahedral angle~\cite{stillinger1985computer}; Tersoff-type potentials let the strength of a bond depend on its local coordination. With some ten parameters fitted to a few properties, these classical potentials reproduce the crystal well and the configurations far from it less well. For the defect energies and planar faults of silicon, the analytical potentials differ from density-functional values by tens of percent for at least some of the defects~\cite[Sect.~III\,E]{bartok2018machine}.

A machine-learned potential keeps \cref{eq:E.11-local} but lets $E^{\mathrm{at}}_\kappa$ be a flexible function of the neighborhood, fitted to density-functional data. The neighborhood is first encoded in a descriptor that is invariant under rotations, reflections and permutations of like atoms. Examples are symmetry functions of the neighbor distances and angles feeding a neural network~\cite{behler2007generalized}, the smooth overlap of atomic positions feeding a kernel regression~\cite{bartok2010gaussian}, or a systematic expansion in body-ordered invariants, the atomic cluster expansion~\cite{drautz2019atomic}. Equivariant message-passing networks build the descriptor from features that rotate with the structure and are passed between neighbors over several layers~\cite{batzner2022e3,batatia2022mace}. The families differ in cost and in how much data they need; they share \cref{eq:E.11-local} and therefore share its failure modes.

The Gaussian approximation potential for silicon~\cite{bartok2018machine} is a well-documented example and the one we use as the silicon case. Its energy is a fixed pair term plus a kernel sum,
\begin{equation}
  E^{\mathrm{ad}} = \sum_{I<J}V^{(2)}(|\vect{R}_I - \vect{R}_J|) + \sum_I\sum_{s=1}^{M_{\mathrm{rep}}}\alpha_s\,K_{\mathrm{GAP}}(\mathcal{N}_I, \mathcal{N}_s),
  \label{eq:E.11-gap}
\end{equation}
with $K_{\mathrm{GAP}}$ a similarity kernel between the neighborhood of atom $I$ and $M_{\mathrm{rep}}$ representative neighborhoods $\mathcal{N}_s$ drawn from the training data, and coefficients $\alpha_s$ fitted by regularized least squares~\cite[Eq.~(1)]{bartok2018machine}. The pair term $V^{(2)}$ is purely repulsive and is fitted to the density-functional energy of an isolated pair of silicon atoms; it describes the steep exchange repulsion at short distance and leaves the bonding to the kernel term~\cite[Sect.~II\,A]{bartok2018machine}. The cutoff radius $r_{\mathrm{c}}$ is \SI{5}{\angstrom}.

\subsection{Training sets}

Each density-functional calculation in the training set yields a total energy, six components of the stress and three force components per atom, and all of them enter the fit with tolerances that set their relative weight. For silicon the database holds 2475 periodic cells with 171\,815 atomic environments~\cite[Table~I]{bartok2018machine}. The cells cover the diamond crystal at strains and temperatures up to melting, other crystal phases ($\beta$-tin, simple hexagonal, bc8, fcc, hcp and others), liquid and amorphous cells, low-index surfaces with their reconstructions and decohesion paths, the vacancy, the divacancy and the interstitials in cells of 63 to 217 atoms, crack tips and a screw dislocation. The reference is the PW91 functional, with a plane-wave cutoff of \SI{250}{eV}~\cite[Sect.~II\,B]{bartok2018machine}.

A recoil needs the training set to cover two regions, the defects and the short distances. The defects are there by design. The short distances are not, and thermal sampling cannot put them there. At \SI{1.0}{\angstrom} two silicon atoms repel each other with about \SI{51}{eV} (\cref{eq:E.11-zbl} below), while the hottest training configurations, liquid silicon near \SI{2000}{K}, have $T \approx \SI{0.17}{eV}$. A thermal ensemble visits a pair separation of \SI{1.0}{\angstrom} with a Boltzmann weight of order $\ee^{-51/0.17}$, which is zero for every practical purpose. A training set assembled from equilibrium and thermal configurations therefore leaves the potential unconstrained exactly where a recoil cascade begins.

\subsection{Close collisions and the Ziegler--Biersack--Littmark splice}

At separations well below a bond length the energy of two atoms is dominated by the Coulomb repulsion of the nuclei, screened by the electrons that remain between them. The standard description is the universal screened potential of Ziegler, Biersack and Littmark (ZBL)~\cite[Chap.~2]{ziegler1985stopping},
\begin{equation}
  V_{\mathrm{ZBL}}(r) = \frac{Z_1 Z_2\,\alpha}{r}\,\phi\!\left(\frac{r}{a_{\mathrm{U}}}\right), \qquad
  a_{\mathrm{U}} = \frac{0.8854\,a_0}{Z_1^{0.23} + Z_2^{0.23}},
  \label{eq:E.11-zbl}
\end{equation}
with $\phi(x) = 0.1818\,\ee^{-3.2x} + 0.5099\,\ee^{-0.9423x} + 0.2802\,\ee^{-0.4029x} + 0.02817\,\ee^{-0.2016x}$, $a_0$ the Bohr radius and $Z_{1,2}$ the nuclear charges. The ZBL potential is a fit to the repulsion of many atom pairs, computed from their free-atom charge densities, which makes it universal in the sense that one screening function serves all pairs once lengths are measured in $a_{\mathrm{U}}$. For silicon on silicon, $a_{\mathrm{U}} = \SI{0.128}{\angstrom}$, and the unscreened Coulomb energy $Z^2\alpha/r = \SI{2822}{eV\,\angstrom}/r$ is reduced by $\phi$ to \SI{51}{eV} at \SI{1.0}{\angstrom} and \SI{9.6}{eV} at \SI{1.5}{\angstrom}.

In a head-on collision of the recoiling atom with a neighbor at rest, the two equal masses share the kinetic energy $E_R/2$ of the center-of-mass frame, and the turning point $r_{\mathrm{min}}$ solves $V_{\mathrm{ZBL}}(r_{\mathrm{min}}) = E_R/2$. Solving this condition numerically with \cref{eq:E.11-zbl} for silicon gives turning points of 1.49, 1.20, 1.00, 0.83, 0.49 and \SI{0.18}{\angstrom} at $E_R = 20$, 50, 100, 200, 1000 and \SI{10000}{eV} (\cref{tab:E-md-si}), against the bond length of \SI{2.35}{\angstrom}. Head-on collisions are the extreme case, and glancing ones stay farther apart. Near the threshold, $E_R \approx 20$ eV, the atoms come to about \SI{1.5}{\angstrom}, inside the bonding region that a good training set covers, while above \SI{100}{eV} they reach the core region that no thermal training set covers.

The two descriptions are joined by a smooth switch,
\begin{equation}
  E^{\mathrm{ad}} = \sum_{I<J}\bigl[1 - f_{\mathrm{sw}}(r_{IJ})\bigr]V_{\mathrm{ZBL}}(r_{IJ}) + f_{\mathrm{sw}}(r_{IJ})\,E^{\mathrm{ad}}_{\mathrm{ML}},
  \qquad f_{\mathrm{sw}}(r) = \frac{1}{1 + \ee^{-b_{\mathrm{sw}}(r - r_{\mathrm{sw}})}},
  \label{eq:E.11-splice}
\end{equation}
written here schematically: in practice the switch multiplies the pair part of the machine-learned energy, and making the many-body part vanish at short distance is left to the training of the potential. The center $r_{\mathrm{sw}}$ and steepness $b_{\mathrm{sw}}$ are chosen so that the potential is pure $V_{\mathrm{ZBL}}$ below about \SI{1}{\angstrom} and pure machine-learned energy above about \SI{1.5}{\angstrom}. The design rule follows from the turning points. Threshold dynamics must not feel the splice, so the machine-learned energy must be fully restored above the $r_{\mathrm{min}}$ of a \SI{20}{eV} collision, \SI{1.49}{\angstrom}. The splice also must not leave a minimum, a shoulder or a jump in the force between the pure $V_{\mathrm{ZBL}}$ region and the pure machine-learned region; a minimum, a shoulder or a jump there would deflect every collision that reaches those separations. The splice is checked against density-functional energies of a silicon pair, and of a crystal with one atom pushed toward a neighbor, at separations between 0.8 and \SI{1.8}{\angstrom}. When the pair term of the potential was already fitted to the density-functional dimer, as in \cref{eq:E.11-gap}, the potential is usually close to $V_{\mathrm{ZBL}}$ near \SI{1}{\angstrom}, and the check then confirms the agreement.

\subsection{Integrating a recoil}

The equations of motion are integrated with the velocity Verlet scheme, which is time-reversible and conserves a shadow energy close to the true one, so that the total energy of an isolated cell drifts only through round-off and through forces that are not smooth. The time step is set by the fastest atom. A silicon atom with \SI{100}{eV} moves at $v = (2E_R/M_\kappa)^{1/2} = \SI{26.2}{km/s}$, which is \SI{0.26}{\angstrom} per femtosecond; at \SI{1}{keV} it moves \SI{0.83}{\angstrom/fs}. To resolve a collision whose turning point lies within a fraction of an \si{\angstrom} of the steep part of $V_{\mathrm{ZBL}}$, the integrator needs steps in which no atom moves more than a few hundredths of an \si{\angstrom}. Moving \SI{0.05}{\angstrom} per step means \SI{0.19}{fs} at \SI{100}{eV} and \SI{0.06}{fs} at \SI{1}{keV}. After a few hundred femtoseconds the fastest atom carries a few eV, and the step can grow to the \SI{1}{fs} that thermal silicon allows, about a sixty-fifth of the shortest phonon period (\SI{65}{fs}, from the \SI{64.1}{meV} optical phonon). Recoil codes therefore adapt the step to the current maximum velocity and force, a practice the silicon and germanium threshold studies follow~\cite[Table~I]{sassi2022energy}. The check on the integration is energy conservation: with the thermostat switched off, the total energy of the cell must stay constant to well below the stored energies being measured, here a small fraction of an eV over the whole run.

The cell has to absorb the recoil. The part of $E_R$ that is not stored in defects spreads over the cell as heat, and the heat must leave through a boundary layer held at the base temperature by a thermostat. The silicon and germanium studies use cells of $8\times8\times8$ conventional cubes (4096 atoms), a thermostatted border \SI{6}{\angstrom} thick, and runs of \SI{20}{ps}~\cite[Sect.~II\,A and Table~I]{sassi2022energy}. The recoil must stop well inside the cell, so that it neither crosses the periodic boundary nor reaches the thermostat while still energetic; at \SI{100}{eV} the primary knock-on atom travels a few \si{\angstrom} to about a nanometer, and $8a = \SI{4.3}{nm}$ is ample. The heat must be removed without reflecting back from the thermostat as a coherent pulse; a border that is too thin or too stiff reflects the heat back as a coherent pulse. And the run must last until the potential energy of the cell has settled; the silicon and germanium studies fixed their run lengths from test runs longer than \SI{5}{ps}, by the time the potential energy took to come within \SI{1}{eV} of its final value~\cite[Sect.~II\,A]{sassi2022energy}.

\subsection{Cascade statistics}

A single trajectory answers a yes-or-no question for one direction and one energy. The physical quantities are averages over the direction of the recoil, which a dark-matter or neutron scattering event draws from its own angular distribution, and over the position of the struck atom within the unit cell. The standard protocol draws $N_{\mathrm{dir}}$ directions uniformly on the sphere and, for each, increases $E_R$ in steps of about \SI{1}{eV} through the threshold region and more coarsely above it. The published silicon and germanium scans use more than 1000 directions~\cite[Sect.~II\,A]{sassi2022energy}.

Whether a defect survived is decided by an occupation count: each atom of the final configuration is assigned to the Wigner--Seitz cell of the nearest perfect-lattice site, and any site with occupation other than one marks a vacancy or an interstitial. The method has no free parameter and does not count the bond-switching defects that relax away~\cite[Sect.~II\,A]{sassi2022energy}. From the counts we obtain $p_{\mathrm{FP}}(E_R)$, the fraction of directions with a surviving defect, with the binomial error $[p_{\mathrm{FP}}(1 - p_{\mathrm{FP}})/N_{\mathrm{dir}}]^{1/2}$, at most \num{0.016} for 1000 directions. Along many directions the defect appears at one energy, disappears at a slightly higher energy because the interstitial now lands close enough to its vacancy to recombine, and reappears at still higher energy. Silicon and germanium show some of this recombination; sapphire shows so much that a direction-resolved threshold is ill-defined~\cite[Sect.~II\,B]{sassi2022energy}. The rate formulas of \cref{sec:19.4} therefore need $p_{\mathrm{FP}}(E_R, \hat{\vect{r}}_R)$ itself, and $E_{\mathrm{d}}(\hat{\vect{r}}_R)$ is quoted as the lowest recoil energy at which $p_{\mathrm{FP}}$ reaches one half (\cref{eq:19.1-Ed}). With the Stillinger--Weber potential the minimum threshold over all directions is \SI{18}{eV} in silicon and \SI{13}{eV} in germanium (with a modified Stillinger--Weber potential); above it the defect probability rises steeply, though more slowly than in diamond, because of recombination~\cite[Table~II, Fig.~1 and Sect.~II\,B]{sassi2022energy}. The silicon and germanium minima are both properties of the potentials, and the germanium value is model-conditional. Kadribasic and coworkers give 17.5 and \SI{12.5}{eV} for the same two potentials, the minima of threshold surfaces averaged over eight runs~\cite{kadribasic2018directional}.

\subsection{Stored energy and the quench protocol}

The stored energy of a trajectory is the potential energy the cell keeps once all motion has died out. The published scans take the difference of potential energies at the end and at the start of a finite run~\cite[Sect.~II\,A]{sassi2022energy}, and this difference includes whatever heat has not yet left the cell. The error from the residual heat follows from equipartition: in the harmonic crystal the potential energy is $\tfrac32N_{\mathrm{at}}T$, so a residual temperature rise of \SI{1}{K} in 4096 atoms adds $\tfrac32\times4096\times\SI{8.6e-5}{eV} = \SI{0.53}{eV}$, comparable to the uncertainty of the stored energy itself. A quench removes the residual heat. The final configuration is relaxed to its local minimum by energy minimization: conjugate gradients, or damped dynamics with the velocities set to zero whenever the power $-\nabla E^{\mathrm{ad}}\cdot\dot{\vect{R}}$ turns negative. The perfect crystal is relaxed the same way, and
\begin{equation}
  E_{\mathrm{st}} = E^{\mathrm{ad}}_{\mathrm{min}}(\text{final}) - E^{\mathrm{ad}}_{\mathrm{min}}(\text{perfect}) .
  \label{eq:E.11-Est}
\end{equation}
The minimum reached is the inherent structure of the final configuration: thermal motion is removed and the defects are kept, as long as the barriers between them exceed the small forces of the minimizer. For the experiments of \cref{sec:19.4}, which measure heat at millikelvin long after the event, \cref{eq:E.11-Est} is the relevant quantity.

A second protocol choice concerns the reference energies. A cascade census produces vacancies, interstitials of several types, close Frenkel pairs and small clusters, each with a formation energy that the potential predicts with some error. The stored energy of a trajectory is roughly the sum of its defects' energies plus their interactions. One table of configurational energies of the recurring defects is computed once with density-functional theory in large relaxed cells, and every census is checked against it. The table makes results from different potentials, or from a potential and the density functional, comparable. For the silicon potential of \cref{eq:E.11-gap}, the formation energies of the vacancy and of the hexagonal, tetrahedral and dumbbell interstitials lie within \SI{7}{\percent} of the density-functional reference~\cite[Sect.~III\,E]{bartok2018machine}. A defect storing several eV then carries an error of a few tenths of an eV from the potential alone, which is the floor on how precisely a single-pair stored energy can be quoted.

\subsection{Defect paths and barriers}

Once defects survive, \cref{sec:19.6} needs to know how they move and recombine. The tool is the nudged elastic band: a chain of $N_{\mathrm{im}}$ replicas of the cell, between two relaxed end states, is relaxed with spring forces along the chain and the true force perpendicular to it, until the chain lies on the minimum-energy path. In the climbing-image variant the highest replica is driven uphill along the path and converges onto the saddle point, whose energy relative to the initial minimum is the barrier $V_{\mathrm{b}}$~\cite{henkelman2000climbing}. End states come from the cascade census, which at $E_R = \SI{100}{eV}$ already supplies close pairs at many separations and orientations. Paths are first mapped with the machine-learned potential, at a cost of a few thousand force calls each, and the saddle points that matter are recomputed with density-functional theory in cells of 64 to 216 atoms. The silicon potential reproduces the density-functional ordering of the hexagonal and tetrahedral interstitials, which are related by a short displacement along a diffusion path, though the potential underestimates their energy difference; the analytical potentials make larger errors, and many reverse the order~\cite[Sect.~III\,E]{bartok2018machine}.

A barrier becomes a rate through harmonic transition-state theory, $\varGamma = \varGamma_0\,\ee^{-V_{\mathrm{b}}/T}$, with an attempt rate $\varGamma_0$ from the vibrational frequencies at the minimum and at the saddle. At the operating temperature of a cryogenic detector the transition-state rate is useless for any barrier of chemical size: at \SI{10}{mK}, $T = \SI{0.86}{\micro eV}$, and a barrier of \SI{0.1}{eV} gives a Boltzmann factor of $\ee^{-1.2\times10^5}$. At millikelvin, defects move by tunneling through low barriers, or are activated by the athermal phonons of a later event. Tunneling needs the barrier shape and the tunneling mass along the path, which the nudged elastic band provides and \cref{sec:19.6} turns into a rate; activation by phonons needs the energy that reaches the defect, which the transport of \cref{app:E.12} provides.

\subsection{Convergence parameters and cost}

For the potential: the size and coverage of the training set, the cutoff $r_{\mathrm{c}}$, the descriptor resolution, and the splice $(r_{\mathrm{sw}}, b_{\mathrm{sw}})$. For the dynamics: the time-step criterion, the cell size, the border thickness and the run length. For the statistics: $N_{\mathrm{dir}}$, the energy grid, and the sampling of the struck atom over the inequivalent sites of the cell. For the stored energy: the quench tolerance. Each is converged by varying it and checking that $p_{\mathrm{FP}}$ and the mean $E_{\mathrm{st}}$ move by less than their statistical errors.

The cost is set by the number of force evaluations. The silicon Gaussian approximation potential costs about \SI{100}{ms} per atom per force call on one core~\cite[Sect.~V]{bartok2018machine}. A 4096-atom cell then needs \SI{410}{s} per step, and a trajectory of $2\times10^4$ steps ($\sim\SI{20}{ps}$ at a typical \SI{1}{fs} step after the collision phase) needs $8\times10^6$ core-seconds, about $2\times10^3$ core-hours. A scan of 1000 directions at 20 energies is $2\times10^4$ trajectories, or $5\times10^7$ core-hours. A production study with a machine-learned potential therefore has to use one of the faster descriptor families, on graphics processors, and has to spend its trajectories where $p_{\mathrm{FP}}$ changes fastest. A density-functional reference trajectory in a 216-atom cell costs far more per atom, but only a few directions are needed to validate the potential.

\subsection{Pitfalls}

The most common failure is a potential that is accurate at equilibrium and wrong in the region a recoil explores (\cref{sec:19.1-md}). The remedy is to validate the threshold itself: thresholds along a few directions from density-functional molecular dynamics, the approach of the density-functional reference study for silicon~\cite{holmstrom2008threshold}.

\subsection{Silicon}

\Cref{tab:E-md-si} collects the silicon numbers of this section.

\begin{table}[t]
  \caption{Recoil molecular dynamics in silicon: the inputs and the quantities a production study converges.}
  \label{tab:E-md-si}
  \centering
  \begin{tabular}{p{0.40\linewidth}p{0.48\linewidth}}
    \toprule
    Quantity & Silicon \\
    \midrule
    Bond length; shortest phonon period & \SI{2.35}{\angstrom}; \SI{65}{fs} \\
    Head-on turning point $r_{\mathrm{min}}$ at $E_R = 20$, 50, 100, 200, 1000, \SI{10000}{eV} & 1.49, 1.20, 1.00, 0.83, 0.49, \SI{0.18}{\angstrom} (\cref{eq:E.11-zbl}) \\
    Speed of the struck atom at 100 and \SI{1000}{eV} & 0.26 and \SI{0.83}{\angstrom/fs} \\
    Training set of the machine-learned potential & 2475 cells, 171\,815 environments, PW91~\cite{bartok2018machine} \\
    Point-defect formation energies vs density-functional theory & within \SI{7}{\percent}~\cite{bartok2018machine} \\
    Cost per atom per force call & \SI{100}{ms} on one core~\cite{bartok2018machine} \\
    Cell, border, run length & 4096 atoms, \SI{6}{\angstrom}, \SI{20}{ps}~\cite{sassi2022energy} \\
    Directions sampled & $> 1000$~\cite{sassi2022energy} \\
    Minimum threshold (Stillinger--Weber) & \SI{18}{eV}~\cite{sassi2022energy} \\
    Residual heat per K in 4096 atoms & \SI{0.53}{eV} (equipartition) \\
    \bottomrule
  \end{tabular}
\end{table}
